\documentclass[11pt, a4paper]{article}

% --- PACCHETTI ---
\usepackage[utf8]{inputenc}
\usepackage[italian]{babel}
\usepackage[margin=2.5cm]{geometry} % Imposta i margini della pagina
\usepackage{graphicx}               % Per inserire il logo
\usepackage{tabularx}               % Per tabelle larghe quanto la pagina
\usepackage[table]{xcolor}          % Per colorare le celle delle tabelle
\usepackage[hidelinks]{hyperref}    % Per stile link cliccabili
\usepackage{lmodern}                % Rende i font scalabili all'infinito
\usepackage{anyfontsize}            % Permette di usare dimensioni personalizzate
\usepackage[T1]{fontenc}

\begin{document}

% ==============================================================================
% FRONTESPIZIO (PAGINA 1)
% ==============================================================================
\begin{titlepage}
    \centering
    \vspace*{0.5cm}

    % LOGHI
    \includegraphics[width=0.45\textwidth]{assets/logo_unipd.jpg}
    \hfill
    \includegraphics[width=0.45\textwidth]{assets/logo.png}\par
    \vspace{0.5cm}

    % NOME DEL GRUPPO
    {\fontsize{45}{50}\sffamily\selectfont \textbf{On-By-One}\par} 
    \vspace{1cm}

    % TITOLO DEL VERBALE
    {\huge \bfseries Verbale interno 2026-11-23 \par}
    
    \vspace{1.5cm}

    % MEMBRI DEL GRUPPO
    {\large
    \textbf{On-By-One}:\par
    \vspace{0.3cm}
    \begin{tabular}{l l}
        Sara Ballardin & 2101048 \\
        Shuhui Chen & 2138001 \\
        Simone Giuliano & 2145591 \\
        Alvise Greggio & 2141004 \\
        Daniele Nuzzaci & 2139039 \\
        Federico Temporin & 2146764 \\
        Tommaso Zorzetto & 2148629 \\
    \end{tabular}
    \par}

    \vfill

    % TABELLA INFORMAZIONI DOCUMENTO E DESTINATARI
    \begin{table}[h]
        \centering
        \arrayrulecolor{gray!50}
        \renewcommand{\arraystretch}{1.5}
        \begin{tabularx}{\textwidth}{|c|c|c|X|}
            \hline
            \rowcolor{gray!10} \multicolumn{4}{|c|}{\textbf{Informazioni documento}} \\
            \hline
            \rowcolor{gray!10} \textbf{Versione} & \textbf{Data} & \textbf{Stato} & \multicolumn{1}{c|}{\textbf{Destinatari}} \\
            \hline
            - & - & Bozza & 
            On-By-One \newline 
            Prof. Tullio Vardanega \newline 
            Prof. Riccardo Cardin \\
            \hline
        \end{tabularx}
        \arrayrulecolor{black}
    \end{table}
    
    \vspace{0.5cm}
    
    % CONTATTI
    {\large Contatti: \href{mailto:onbyone.swe@gmail.com}{onbyone.swe@gmail.com} }
    
\end{titlepage}

\newpage

% ==============================================================================
% REGISTRO DELLE MODIFICHE (PAGINA 2)
% ==============================================================================
\section*{Registro delle modifiche}

\vspace{1cm}
\begin{center}
    \arrayrulecolor{gray!50}
    \renewcommand{\arraystretch}{1.5}
    \begin{tabularx}{\textwidth}{|c|c|c|X|c|}
        \hline
        \rowcolor{gray!10} 
        \multicolumn{1}{|c|}{\textbf{Versione}} & 
        \multicolumn{1}{c|}{\textbf{Data}} & 
        \multicolumn{1}{c|}{\textbf{Autore}} & 
        \multicolumn{1}{c|}{\textbf{Descrizione}} & 
        \multicolumn{1}{c|}{\textbf{Verificatore}} \\
        \hline
        0.2.0 & 2026-11-24 & Sara Ballardin & Stesura del verbale interno & - \\
        \hline
        0.1.0 & 2026-11-23 & Sara Ballardin & Bozza iniziale & - \\
        \hline
    \end{tabularx}
    \arrayrulecolor{black}
\end{center}
\vfill
\newpage

% ==============================================================================
% INDICE (PAGINA 3)
% ==============================================================================
\tableofcontents
\newpage

% ==============================================================================
% INFORMAZIONI GENERALI
% ==============================================================================
\section{Informazioni generali}

\begin{table}[h]
    \centering
    \arrayrulecolor{gray!50}
    \renewcommand{\arraystretch}{1.5}
    \begin{tabularx}{\textwidth}{|c|>{\centering\arraybackslash}X|}
        \hline
        \rowcolor{gray!10} \multicolumn{2}{|c|}{\textbf{Dettagli della riunione}} \\
        \hline
        \textbf{Data} & 2026-11-23 \\ \hline
        \textbf{Ora inizio} & 15:00 \\ \hline
        \textbf{Ora fine} & 17:50 \\ \hline
        \textbf{Modalità} & Discord \\ \hline
    \end{tabularx}
    \arrayrulecolor{black}
\end{table}

\begin{table}[h]
    \centering
    \arrayrulecolor{gray!50}
    \renewcommand{\arraystretch}{1.5}
    \begin{tabularx}{\textwidth}{|>{\centering\arraybackslash}X|>{\centering\arraybackslash}X|>{\centering\arraybackslash}X|}
        \hline
        \rowcolor{gray!10} \multicolumn{3}{|c|}{\textbf{Partecipanti}} \\
        \hline
        \rowcolor{gray!10} \textbf{Nome e cognome} & \textbf{Ruolo} & \textbf{Presenza} \\
        \hline
        Sara Ballardin & Verificatore & Presente \\ \hline
        Shuhui Chen & Responsabile & Presente \\ \hline
        Simone Giuliano & Amministratore & Presente \\ \hline
        Alvise Greggio & Analista & Presente \\ \hline
        Daniele Nuzzaci & Amministratore & Presente \\ \hline
        Federico Temporin & Analista & Presente \\ \hline
        Tommaso Zorzetto & Analista & Presente \\ \hline
    \end{tabularx}
    \arrayrulecolor{black}
\end{table}

% ==============================================================================
% CORPO DEL VERBALE
% ==============================================================================

\section{Ordine del giorno}
\begin{itemize}
    \item Punto sui lavori svolti, difficoltà riscontrate e retrospettiva dello sprint concluso.
    \item Rotazione e ribilanciamento dei ruoli per lo sprint 2.
    \item Revisione e organizzazione dell'esposizione per l'incontro SAL con Vimar S.p.A. del 2026-11-24.
    \item Valutazione degli strumenti UML, gestione delle \textit{milestone} e revisione dell'automazione documentale.
\end{itemize}

\section{Discussioni}
La riunione si è aperta con un confronto sullo stato dei lavori svolti e riorganizzando i ruoli per lo sprint 2. Considerato il carico legato all'infrastruttura e alla gestione dei \textit{workflow}, il gruppo ha deciso di ribilanciare le risorse spostando un membro dal ruolo di Verificatore a quello di Amministratore. Subito dopo si è svolta una retrospettiva sullo sprint appena concluso per confrontarsi sulle difficoltà incontrate e raccogliere consigli pratici a supporto di chi subentra nei nuovi incarichi.

\noindent Sul piano organizzativo, il team ha rivisto le slide già predisposte per l'incontro SAL con Vimar S.p.A. del 2026-11-24, confermandone i contenuti e verificando che non ci fossero altri dubbi da aggiungere; ci si è poi suddivisi l'esposizione dei vari punti in base ai ruoli ricoperti durante lo sprint. Per i diagrammi UML, vista l'incertezza tra \texttt{draw.io}, \texttt{Mermaid} e \texttt{PlantUML}, si è scelto di sperimentarli prima per poi confrontarsi su quale strumento risulti più comodo e adatto ai diagrammi da produrre. Per tracciare con più ordine le parti dei documenti che proseguono oltre il singolo sprint, si è deciso di ricorrere alle milestone su GitHub; al momento l'unica individuata è \textit{Casi d'uso - RTB}.

\noindent Infine, sono stati esaminati alcuni problemi dello script che aggiorna in automatico il registro delle modifiche: manca la riga separatrice \texttt{\textbackslash hline}, il campo descrizione va popolato dalla descrizione della \textit{Pull Request} anziché dal titolo (evitando codice LaTeX nei titoli) e va rivista la gestione del versionamento (\texttt{X.Y.Z}). In attesa delle modifiche, l'automazione resta sospesa e si procede manualmente. Si è stabilito di uniformare tutte le date dei documenti nel formato con trattini (\texttt{YYYY-MM-DD}).

\section{Retrospettiva}
Dal confronto sullo sprint appena concluso sono emersi alcuni consigli pratici per chi subentra nei vari ruoli:
\begin{itemize}
    \item \textbf{Responsabile:} Leggere regolarmente tutti i verbali per monitorare rischi e pianificazione, compilare con cura la rendicontazione delle ore e tenere aggiornato passo dopo passo il \textit{Piano di Progetto}. Sfruttare le tabelle dei \textit{to-do} sia per il \textit{Diario di Bordo} sia per i SAL aziendali, favorendo un confronto aperto con il gruppo per evitare fraintendimenti.
    \item \textbf{Amministratore:} Valutare il supporto di strumenti AI con modalità di pianificazione (esempio:\ Copilot, Antigravity), sfruttando i piani studenti per gestire i workflow con una visione d'insieme. È inoltre fondamentale avvisare tempestivamente il team ogni volta che vengono introdotte modifiche agli strumenti o alle modalità operative.
    \item \textbf{Analista:} Iniziare da una visione d'insieme per suddividere il lavoro in modo equilibrato, partendo sempre dallo studio delle tecnologie coinvolte per comprenderne il funzionamento reale prima della stesura dei casi d'uso.
    \item \textbf{Verificatore:} Registrare le riunioni o prendere appunti precisi e tenere sempre aggiornati rendicontazione, tabelle delle PR verificate e \textit{issue} in stato \textit{draft}, così da non accumulare arretrati nelle revisioni.
\end{itemize}

\section{Decisioni}
\begin{itemize}
    \item \textbf{VI-2026-11-23.1} Ribilanciare i ruoli per lo sprint 2 spostando un membro dal ruolo di Verificatore a quello di Amministratore.
    \item \textbf{VI-2026-11-23.2} Sperimentare \texttt{draw.io}, \texttt{Mermaid} e \texttt{PlantUML} prima di scegliere lo strumento definitivo per i diagrammi UML.
    \item \textbf{VI-2026-11-23.3} Utilizzare le milestone su GitHub per tracciare le parti dei documenti che proseguono oltre il singolo sprint, introducendo \textit{Casi d'uso - RTB}.
    \item \textbf{VI-2026-11-23.4} Uniformare tutte le date dei documenti nel formato con trattini (\texttt{YYYY-MM-DD}).
    \item \textbf{VI-2026-11-23.5} Sospendere temporaneamente l'automazione del registro delle modifiche e tornare alla compilazione manuale finché non verranno risolti i problemi riscontrati.
\end{itemize}

\section{To-do}
\begin{center}
    \arrayrulecolor{gray!50}
    \renewcommand{\arraystretch}{1.5}
    \begin{tabularx}{\textwidth}{|c|X|c|c|}
        \hline
        \rowcolor{gray!10} 
        \multicolumn{1}{|c|}{\textbf{ID issue}} & 
        \multicolumn{1}{c|}{\textbf{Task}} & 
        \multicolumn{1}{c|}{\textbf{Assegnatario}} & 
        \multicolumn{1}{c|}{\textbf{Scadenza}} \\
        \hline
         \href{https://github.com/On-By-One-Team/On-By-One.github.io/issues/119}{\#119} & Stesura verbale interno 2026-11-23 & Sara Ballardin & 2026-11-26 \\
        \hline
        \href{https://github.com/On-By-One-Team/On-By-One.github.io/issues/105}{\#105} & Contributo casi d'uso 1 - sprint 2 & Federico Temporin & 2026-12-09 \\
        \hline
        \href{https://github.com/On-By-One-Team/On-By-One.github.io/issues/117}{\#117} & Contributo casi d'uso 2 - sprint 2 & Alvise Greggio & 2026-12-09 \\
        \hline
        \href{https://github.com/On-By-One-Team/On-By-One.github.io/issues/118}{\#118} & Contributo casi d'uso 3 - sprint 2 & Tommaso Zorzetto & 2026-12-09 \\
        \hline
        \href{https://github.com/On-By-One-Team/On-By-One.github.io/issues/110}{\#110} & Modifica automazione tabelle tracciabilità documenti (aggiunta \texttt{\textbackslash hline}, valorizzazione descrizione e versione) & Daniele Nuzzaci o Simone Giuliano & 2026-11-30 \\
        \hline
        \href{https://github.com/On-By-One-Team/On-By-One.github.io/issues/114}{\#114} & Modifica \textit{Norme di progetto} (documentazione) per coerenza con l'automazione delle tabelle di tracciabilità & Daniele Nuzzaci o Simone Giuliano & 2026-11-30 \\
        \hline
        \href{https://github.com/On-By-One-Team/On-By-One.github.io/issues/122}{\#122} & Aggiornamento formato date con trattini (\texttt{YYYY-MM-DD}) nei documenti & Daniele Nuzzaci o Simone Giuliano & 2026-11-30 \\
        \hline
    \end{tabularx}
    \arrayrulecolor{black}
\end{center}

\end{document}